\usepackage[scr=rsfs,cal=euler]{mathalfa}
\newenvironment{enumeratea}
{\bgroup\def\theenumi{\alph{enumi}}\begin{enumerate}}
{\end{enumerate}\egroup}

\let\epsilon\varepsilon
\newcommand\mto{\mathchoice{\longmapsto}{\mapsto}{\mapsto}{\mapsto}}
\newcommand\quand{\quad\text{and}\quad}
\newcommand\qquand{\qquad\text{and}\qquad}
\newcommand{\spfrac}[2]{\sfrac{#1}{(#2)}}

\usepackage{tikz}
\usetikzlibrary{cd}
\tikzcdset{arrow style=Latin Modern}
\newcommand{\dradpl}{{\hspace*{-2.5mm}\begin{tikzcd}[ampersand replacement=\&,column sep = scriptsize]\mbox{}\ar[r,dashed]\&\mbox{}\end{tikzcd}\hspace*{-2.5mm}}}
\newcommand{\dratst}{{\hspace*{-2.5mm}\begin{tikzcd}[ampersand replacement=\&,column sep = small]\mbox{}\ar[r,dashed]\&\mbox{}\end{tikzcd}\hspace*{-2mm}}}
\let\drascst\dratst
\newcommand{\dra}{\mathchoice{\dradpl}{\dratst}{\drascst}{\drascst}}

\usepackage{stmaryrd}

\usepackage[all,cmtip]{xy}
\xyoption{matrix}\let\labelstyle\textstyle

\usepackage{mathtools}

\newtheorem{mytheorem}{Theorem}
\renewcommand*{\themytheorem}{\Alph{mytheorem}}
\newtheorem{lemma}{Lemma}[section]
\newtheorem{corollary}[lemma]{Corollary}
\newtheorem{theorem}[lemma]{Theorem}
\newtheorem{proposition}[lemma]{Proposition}
\newtheorem{question}[lemma]{Question}

\newtheorem{maintheorem}{Theorem}
\renewcommand\themaintheorem{\Alph{maintheorem}}
\newtheorem{maincorollary}[maintheorem]{Corollary}
\newtheorem{mainprop}[maintheorem]{Proposition}

\theoremstyle{definition}
\newtheorem{definition}[lemma]{Definition}
\newtheorem{notation}[lemma]{Notation}

\newtheorem{remark}[lemma]{Remark}
\newtheorem{example}[lemma]{Example}

\newcommand\iso{\stackrel{\simeq}{\longrightarrow}}

\newcommand\A{{\mathbb A}}\let\AA\A
\newcommand\C{{\mathbb C}}\let\CC\C
\renewcommand\k{{\mathbf k}}
\newcommand\p{{\mathbb P}}\let\PP\p
\newcommand\Z{{\mathbb Z}}\let\ZZ\Z
\newcommand\FF{{\mathbb F}}
\newcommand\GG{{\mathbb G}}
\newcommand\NN{{\mathbb N}}
\newcommand\QQ{{\mathbb Q}}

\newcommand{\kk}{\mathbf{k}}

\let\rat\dra

\newcommand\mm{{\mathfrak m}}

\DeclareMathOperator{\Aut}{Aut}
\DeclareMathOperator{\car}{char}
\DeclareMathOperator{\sing}{sing}
\DeclareMathOperator{\reg}{reg}
\DeclareMathOperator{\Pic}{Pic}
\DeclareMathOperator{\Spec}{Spec}
\DeclareMathOperator{\Var}{Var}
\newcommand{\et}{\textup{ét}}

\newcommand{\set}[2]{\{#1 \mid #2\}}
\newcommand{\Bigset}[2]{\Bigl\{#1 \Bigm| #2\Bigr\}}

\hyphenation{co-cycle}
